\section{Introduction}

In this chapter we will look at the connection between the geometry of
ellipses and the physical constants that we've shown previously.

\section{The Apses}

The apses are the closest and furthest point from the focus in the orbit. 
In other words they are the maximum and minimum radius. At the apses
\begin{equation}
\frac{d r}{dt}=\dot{r} = 0
\end{equation}
We can then use equation \refpage{eq:gravity6} to show that the velocity is
perpendicular to the radius at the apses.
\begin{align}
\dot{r} &= \frac{\mathbf{r}\cdot\mathbf{v}}{r} \nonumber \\
0 &= \frac{\mathbf{r}\cdot\mathbf{v}}{r} \nonumber \\
0 &= \mathbf{r}\cdot\mathbf{v}
\end{align}
since the scalar product is zero, the two vectors must be 
perpendicular\footnote{or zero, but we are looking at non-zero solutions in
particular}. 

\subsection{Angular Momentum}

If the radius and velocity are perpendicular at the apses, then the magnitude
of the specific angular momentum is
\begin{equation}
h=p v_p = a v_a
\end{equation} 
where $v_p$ is the velocity at the periapsis and $v_a$ is the velocity at apoapsis.

\subsection{Energy}

We can now look at the specific total energy again, firstly at periapsis
\begin{align}
\epsilon &= \frac{v_p^2}{2}-\frac{\mu}{p} \nonumber \\
&= \frac{h^2}{2p^2}-\frac{\mu}{p} \label{eq:phys-geom-energy1}
\end{align} 
then at apoapsis
\begin{align}
\epsilon &= \frac{v_a^2}{2}-\frac{\mu}{a} \nonumber \\
&= \frac{h^2}{2a^2}-\frac{\mu}{a}
\end{align} 
Since total energy is a constant these must be equal
\begin{align}
\frac{h^2}{2p^2}-\frac{\mu}{p} &= \frac{h^2}{2a^2}-\frac{\mu}{a} \nonumber \\
\frac{h^2}{2p^2}-\frac{h^2}{2a^2} &= \frac{\mu}{p} -\frac{\mu}{a} \nonumber \\
\frac{a^2h^2}{2a^2p^2}-\frac{p^2h^2}{2a^2p^2} &= \frac{a\mu}{ap} -\frac{p\mu}{ap} \nonumber \\
\frac{a^2h^2-p^2h^2}{2a^2p^2} &= \frac{a\mu-p\mu}{ap} \nonumber \\
h^2 \frac{a^2-p^2}{2 ap} &= (a-p)\mu \nonumber \\
h^2 &= \frac{2ap(a-p)}{a^2-p^2}\mu \nonumber \\
 &= \frac{2ap(a-p)}{(a-p)(a+p)}\mu \nonumber \\
&= \frac{2ap}{a+p}\mu \label{eq:phys-geom1}
\end{align}

\subsection{Velocity}
We now have an equation\refpage{eq:phys-geom1} for the specific angular momentum based purely
on the central mass and the geometry of the orbit. We can use this to
calculate velocities at the periapsis
\begin{equation}
v_p = \frac{h}{p} = \frac{\sqrt{\mu}}{p}\sqrt{\frac{2ap}{a+p}}
\end{equation}
and for the apoapsis
\begin{equation}
v_a = \frac{h}{a} = \frac{\sqrt{\mu}}{a}\sqrt{\frac{2ap}{a+p}}
\end{equation}

\subsection{Semi-Major Axis}
We can feed this result back into equation \refpage{eq:phys-geom-energy1}
\begin{align}
\epsilon &= \frac{h^2}{2p^2}-\frac{\mu}{p} \nonumber \\
&= \frac{2ap}{2p^2(a+p)}\mu-\frac{\mu}{p} \nonumber \\
&= \frac{\mu}{p}\left(\frac{a}{a+p}-1 \right) \nonumber \\
&= \frac{\mu}{p}\left(\frac{a}{a+p}-\frac{a+p}{a+p} \right) \nonumber \\
&= \frac{\mu}{p}\left(\frac{a-a-p}{a+p}\right) \nonumber \\
&= \frac{\mu}{p}\left(\frac{-p}{a+p}\right) \nonumber \\
&= \frac{-\mu}{a+p}
\end{align}
Using equation \refpage{eq:geom3} we can replace the $a+p$ term
\begin{align}
\epsilon &= \frac{-\mu}{2A} \nonumber \\
A &= \frac{-\mu}{2\epsilon}\label{eq:phys-geom5}
\end{align}

\subsection{Eccentricity}

We will now use the results of equation \refpage{eq:phys-geom1}

\begin{equation*}
h^2= \frac{2ap}{a+p}\mu
\end{equation*}
as before we can use equation \refpage{eq:geom3} to relpace $a+p$, additionally
we can use equation \refpage{eq:geom5} to replace the $ap$ term.
\begin{align}
h^2 &= \frac{2A^2(1-e^2)}{2A}\mu \nonumber \\
&= A(1-e^2)\mu
\end{align}
now we can use equation \refpage{eq:phys-geom5} to replace $A$
\begin{align}
h^2 &= A(1-e^2)\mu \nonumber \\
&=\frac{-\mu}{2\epsilon}(1-e^2)\mu \nonumber \\
\frac{2\epsilon h^2}&=e^2-1\nonumber \\
1+\frac{2\epsilon h^2}{\mu^2}&=e^2
\end{align}